\section{Cote 123: the rational singularities of adjoint-quotient fibres}
\label{sec:SM-cote123-rational}

The letter to Kostant dated 22 October 1969, physical pages 22--23
of cote 123 (handwritten pages 21--22), is explicitly over
$\mathbb C$.  It asks whether the fibres of
$\mathfrak g\to\mathfrak h/W$, and of $G\to T/W$ when $G$ is
simply connected, have rational singularities.  Its stated test is
$R^if_*\mathcal O=0$ for $i>0$ on a desingularization.
This is a known positive result, independent of any adopted
analytic-number-theory input.

\paragraph{Finding F-SM-23 (known affirmative answer, with the reduction written out).}
Let $\mathfrak g$ be a complex semisimple Lie algebra.
Every scheme-theoretic fibre of its adjoint quotient
$\mathfrak g\to\mathfrak h/W$ is reduced and normal and has
rational singularities.  Let $G$ be a connected simply connected
complex semisimple algebraic group, with maximal torus $T$ and Weyl
group $W$.  The same assertions hold for every scheme-theoretic
fibre of $G\to T/W$.
Thus for each such fibre $F$ and every proper resolution
$f:Y\to F$,
\[
 f_*\mathcal O_Y=\mathcal O_F,\qquad
 R^if_*\mathcal O_Y=0\quad(i>0).
\]
These assertions use classical results; the following argument is a
complete deduction from the named inputs, not a new proof of those
inputs.

\begin{proof}
\emph{The Lie-algebra statement.}
Kraft and Procesi (1982, \S0.3, p.~540) state reducedness and
normality for every fibre.  Their immediately following \S0.4 on
the same page explicitly states that these fibres have rational
singularities, citing Hesselink.  The scope of \S0.2 is a reductive
complex group; these introductory assertions are not restricted
to the classical orbit-closure classification treated later in that
paper.  This supplies the first statement directly.

\emph{The group fibre as a scheme.}
Put $r=\operatorname{rank}G$.  Steinberg's fundamental-character
coordinates identify $T/W$ with $\mathbb A^r$ (1965,
Theorem~6.1(b), p.~62, and \S6.16, p.~65).
For a fibre $F$, Theorem~6.11, p.~64, gives an irreducible
underlying variety of codimension $r$, with a dense regular class
whose complement has codimension at least two.
The fibre ideal is generated by $r$ functions in the smooth
variety $G$.  Its height is $r$, so these functions form a regular
sequence locally.  Thus $F$ is Cohen--Macaulay.  Theorem~8.1,
p.~70, says that their differentials are independent on the
regular class.  Therefore $F$ is generically reduced; the absence
of embedded associated primes for a Cohen--Macaulay scheme implies
that $F$ is reduced.  Its complement to the smooth regular locus
has codimension at least two.  Serre's $R_1$ and $S_2$ criterion
then gives normality.  Normality is a deduction here; it is not
misattributed to the literal statement of Theorem~6.11.

\emph{Reduction to a unipotent variety.}
Let $s$ represent the unique closed semisimple class in $F$ and
put $L=Z_G(s)$.  This is a connected reductive group: reductivity
and connectedness are the semisimple-centralizer facts recorded
in Steinberg (1965, \S2.7 and Remark~2.10, p.~53).  The latter
remark states the connectedness result with proof deferred; we
use the established result rather than claim to reproduce that
proof.

Since $F$ is the closure of its regular class, the algebraic
semisimple-part map of Kraft--Procesi (\S0.2, pp.~539--540) gives
a $G$-equivariant morphism
\[
 j:F\longrightarrow G\cdot s=G/L.
\]
Its fibre over $s$ is $s\mathcal U(L)$, where $\mathcal U(L)$ is
the unipotent variety.  Indeed, Jordan decomposition gives that
description for every element with semisimple part $s$; conversely
every $su$, $u\in\mathcal U(L)$, has the same invariant-character
values as $s$ and belongs to $F$.  The standard homogeneous-space
descent, applied to $G\to G/L$, therefore gives
\[
 F \simeq G\times^L \mathcal U(L),\qquad
 [g,u]\longmapsto gsu g^{-1}.
 \tag{SM23.1}
\]
Here the right-hand fibre has its reduced structure.  This also
agrees with the scheme fibre: after the smooth faithfully flat
base change $G\to G/L$, the left side is
$G\times j^{-1}(s)$ and is reduced, since $F$ is reduced.
Faithfully flat descent of reducedness over $\mathbb C$ makes
$j^{-1}(s)$ reduced.  The preceding identification of reduced
closed subvarieties is consequently scheme-theoretic.

Hesselink (1976, Theorem~A and \S2, p.~249) supplies the
equivariant Springer resolution
\[
 \tau:Y_L\longrightarrow\mathcal U(L),\qquad
 \tau_*\mathcal O_{Y_L}=\mathcal O_{\mathcal U(L)},\qquad
 R^i\tau_*\mathcal O_{Y_L}=0\quad(i>0).
 \tag{SM23.2}
\]
The theorem is stated for a split reductive group in
characteristic zero, and hence applies to $L$.
The construction is the Springer resolution of the nilpotent
cone; the equivariant logarithm identifies that cone with
$\mathcal U(L)$ in characteristic zero.  Its equivariance is
essential for the next step.

Form the associated morphism
\[
 f:G\times^L Y_L\longrightarrow G\times^L\mathcal U(L)\simeq F.
\]
The smooth $L$-torsor $G\to G/L$ has sections \'etale locally.
For such a cover $V\to G/L$, the base change of $f$ is
$\operatorname{id}_V\times\tau$.
It is proper and birational, with smooth source, so these
properties descend and $f$ is a resolution.
Flat base change for $\tau$, followed by (SM23.2), gives
\[
 R^i(\operatorname{id}_V\times\tau)_*
       \mathcal O_{V\times Y_L}
 =
 \begin{cases}
 \mathcal O_{V\times\mathcal U(L)},&i=0,\\
 0,&i>0.
 \end{cases}
\]
The corresponding assertions for $Rf_*\mathcal O$ descend along
the resulting faithfully flat \'etale cover of $F$.
This proves rational singularities of $F$.

Finally, the independence of the chosen resolution in
characteristic zero is part of the standard rational-singularity
criterion: the definition with \emph{some} resolution is equivalent
to the definition with \emph{every} resolution.  A directly checked
primary research reference stating this equivalence is
Glazer (2019, \S2A, p.~488), which cites
Kempf--Knudsen--Mumford--Saint-Donat, I.3, pp.~50--51.
It yields the displayed conclusion for every resolution.
\end{proof}

\paragraph{Relation to the particular resolution in the letter.}
Steinberg (1976, Corollary~6.4, p.~223), explicitly attributed
there to Grothendieck, states that for the closure $F$ of a regular
conjugacy class and $t\in T\cap F$ the morphism
\[
 \{(x,gB):g^{-1}xg\in tU\}\longrightarrow F
\]
is a desingularization.  This applies under the present complex
simply connected hypothesis.  F-SM-23 and independence of the
resolution give the requested higher-direct-image vanishing for
this resolution as well.  Corollary~6.4 alone supplies a
resolution, not the higher-direct-image vanishing; the latter
comes from rational singularities as explained above.

\paragraph{Scope and status.}
This closes the rational-singularity question in the original
complex setting.  It is not presented as a new theorem of this
project, as a characteristic-positive analogue, or as a statement
about arbitrary individual orbit closures.  No general
non-simply-connected group assertion is needed here.

\subsection*{References for F-SM-23}
\begin{itemize}
\item Glazer, I. (2019). On rational singularities and counting points
of schemes over finite rings. \emph{Algebra \& Number Theory, 13}(2),
485--500. \url{https://doi.org/10.2140/ant.2019.13.485}.
Exact input: \S2A, p.~488.
\item Hesselink, W. H. (1976). Cohomology and the resolution of the
nilpotent variety. \emph{Mathematische Annalen, 223}(3), 249--252.
\url{https://doi.org/10.1007/BF01360956}.
Exact inputs: Theorem~A and \S2, p.~249.
\item Kraft, H., \& Procesi, C. (1982). On the geometry of conjugacy
classes in classical groups. \emph{Commentarii Mathematici
Helvetici, 57}, 539--602.
\url{https://doi.org/10.1007/BF02565876}.
Exact inputs: \S0.2, pp.~539--540; \S\S0.3--0.4, p.~540.
\item Steinberg, R. (1965). Regular elements of semi-simple algebraic
groups. \emph{Publications Math\'ematiques de l'IH\'ES, 25}, 49--80.
\url{https://doi.org/10.1007/BF02684397}.
Exact inputs: \S2.7 and Remark~2.10, p.~53;
Theorem~6.1(b), p.~62; Theorem~6.11, p.~64;
\S6.16, p.~65; Theorem~8.1, p.~70.
\item Steinberg, R. (1976). On the desingularization of the unipotent
variety. \emph{Inventiones Mathematicae, 36}, 209--224.
\url{https://doi.org/10.1007/BF01390010}.
Exact input: Corollary~6.4, p.~223.
\end{itemize}
